\section{Preparing the openings}\label{sec:preparation}

We first put the fractional openings on a sufficiently fine constant grid.
The use of separated representatives and bundles is motivated by the
discretization approach in~\cite{ByrkaEtAl2026}.  The construction below
rounds a bundle's missing assignment mass together with the facilities
that can supply it.  This coupling provides both an additive opening
bound and a multiplicative connection-cost bound.

Throughout this Section, the client set is nonempty; an instance without
clients requires no cost analysis.  We start with feasible fractional
openings $0\le y_i\le1$.  For each client $j$, choose a nearest-unit
assignment $x_{ij}$, resolving ties in a fixed order, and write
\[
 V_j=\sum_i d(i,j)x_{ij}.
\]
The feasibility constraints are $\sum_i x_{ij}=1$ and $x_{ij}\le y_i$;
there is no supply constraint shared by different clients.

\begin{theorem}\label{thm:preparation}
For every $\eta>0$, there is a constant $\Delta_0$ with the following
property.  Given a finite rational metric instance and feasible rational
openings $0\le y_i\le1$ with a nearest-unit assignment of cost $V_j$
for each client $j$, fix any constant even integer $\Delta\ge\Delta_0$ and let
$g=1/\Delta$.  There is a randomized procedure, running in expected
polynomial time on rational inputs, that produces openings
$\widetilde y_i\in[0,1]\cap g\mathbb Z$, allowing polynomially many
co-located labeled facility pieces, such that
\[
 1\le\sum_i\widetilde y_i\le\sum_i y_i+3g
\]
for every output, and the expected nearest-unit assignment cost of every
client $j$ is at most $(1+\eta)V_j$.
In particular the grid can be required to satisfy any additional
positive instance-independent upper bound on $g$.
\end{theorem}

The expectation in Theorem~\ref{thm:preparation} is taken after rejecting
outcomes that exceed the opening budget.  Section~\ref{sec:derandomization}
gives a deterministic implementation that preserves the corresponding
bound on the sum of the client costs.

We use the following elementary dependent-rounding observation, in both
weighted and unweighted form. The mean-preserving pair update and its
product-supermartingale argument follow the standard
approach~\cite{Srinivasan2001,GandhiEtAl2006}; weighted rounding with one
residual fraction is also studied in~\cite[Section~2]{ByrkaEtAl2017}.
We give the proof to specify exactly the moment guarantees used here.

\begin{lemma}\label{lem:prep-weighted-rounding}
Let $q_e\ge0$ be rational counts and let $\beta_e>0$ be rational weights.
One can randomly update the counts within the bounds
$[\lfloor q_e\rfloor,\lceil q_e\rceil]$, preserving
$\sum_e\beta_e q_e$ on every update, until at most one count is
fractional.  Every first moment is preserved, and the expectation of
every product of two distinct counts is at most the product of their
initial values.  Rounding the possible remaining fractional count with
its correct mean preserves these moment assertions.
For the resulting integer counts $Q_e$,
\[
 \Var(Q_e)\le q_e,\qquad
 \E[Q_eQ_f]\le q_eq_f\quad(e\ne f).
\]
The procedure uses linearly many two-coordinate updates.
\end{lemma}

\begin{proof}
When two coordinates are fractional, move them along the line
\[
 (q_e,q_f)\longmapsto
 (q_e+t,q_f-(\beta_e/\beta_f)t)
\]
to one of the two endpoints allowed by their fixed floor and ceiling
bounds.  Choose the endpoint probabilities so that the conditional
mean of $t$ is zero.  At least one coordinate becomes integral, and
integral coordinates are not updated subsequently.  Each coordinate
is a martingale.  The product of the two updated coordinates is a
concave quadratic in $t$; a product involving exactly one updated
coordinate has unchanged conditional mean.  Thus every distinct-pair
product is a supermartingale.  Rounding the last fractional coordinate
with the correct conditional mean leaves all these products unchanged
in expectation.  Finally, an integer random variable supported on
$\{m,m+1\}$, with mean $m+\theta$, has variance
$\theta(1-\theta)\le m+\theta$.  This proves the variance assertion.
\end{proof}

\begin{proof}[Proof of Theorem~\ref{thm:preparation}]
We first form disjoint bundles containing most of each representative's
assigned mass. We then round their deficits jointly with the openings
that supply them, reserving a small amount of budget for clipping large
deficits. After rounding inside bundles, every representative can again
receive one unit. The final part of the proof compares connection costs
at two representative-distance scales and conditions on the budget event.

We will choose positive rational constants in the order
\[
 \delta,\quad M,\quad T_1,\quad\gamma,\quad g,
\]
with $\delta$ small, $M$ and then $T_1$ large, and $\gamma$ and then $g$
small.  All choices depend only on $\eta$; the last choice means that
every sufficiently fine even grid is permitted.  During the proof we
require
\begin{equation}\label{eq:prep-parameter-basics}
 0<\delta\le\tfrac14,\qquad 0<\gamma\le1,\qquad
 \theta:=\frac1{M\delta}\le\tfrac18,\qquad
 \alpha:=\frac M{\delta T_1}\le\tfrac14.
\end{equation}
Set $p_0=1-\theta$ and $\beta=M/T_1=\delta\alpha$.

\paragraph{Splitting and representatives.}
For each original facility, view its opening $y_i$ as an interval of
length $y_i$ and its assignment to client $j$ as the initial subinterval
of length $x_{ij}$.  Split at the endpoints of these subintervals for all
clients.  There are at most $|J|+1$ positive-mass pieces per original
facility, and every assignment now uses whole pieces.  Retain the
notation $y_i$ for the mass of a piece, and let $F_j$ be the set of pieces
used by client $j$.  Thus
\[
 y(F_j)=1,\qquad
 V_j=\sum_{i\in F_j}y_i d(i,j),\qquad
 y(W):=\sum_{i\in W}y_i.
\]
Pieces retain the distances of their original facility.  They may be
treated as distinct candidates throughout the rounding; merging
co-located integral openings at the end preserves all connection
distances and can only decrease the number of opened facilities.

Process the clients in nondecreasing order of $V_j$.  Select each
remaining client $s$ as a representative and absorb every remaining
client $j$ for which $d(j,s)\le MV_j$.  Write $\mathcal S$ for the set of
representatives.  A client $j$ assigned to $s$ satisfies
\begin{equation}\label{eq:prep-representative-properties}
 V_s\le V_j,\qquad d(j,s)\le MV_j.
\end{equation}
Distinct representatives satisfy
\[
 d(s,t)>M\max(V_s,V_t):
\]
indeed, the later one was not absorbed by the earlier one.  Let $L_s$ be
the distance to a nearest other representative, taking $L_s=\infty$
when $s$ is the only representative.  Define
\[
 P_s=\{i\in F_s:d(s,i)\le\delta L_s\},\quad
 p_s=y(P_s),\quad c_s=1-p_s,\quad A_s=F_s\setminus P_s,
\]
and let $U$ consist of all pieces outside $\bigcup_s P_s$.
For $L_s=\infty$, interpret $P_s=F_s$ and $c_s=0$.

The bundles $P_s$ are disjoint.  Otherwise a piece in $P_s\cap P_t$
would imply
\[
 d(s,t)\le\delta(L_s+L_t)\le2\delta d(s,t),
\]
which is impossible.  For finite $L_s$, the mass in $A_s$ has distance
at least $\delta L_s$ from $s$, so
\begin{equation}\label{eq:prep-deficit-small}
 \delta L_s c_s\le V_s,\qquad
 0\le c_s\le\theta,\qquad p_s\ge p_0.
\end{equation}
These bounds also cover $V_s=0$: in that case $c_s=0$.

\paragraph{Supplier stars and signed counts.}
For the raw opening count, give every bundle a base level of one.  A
supplier star couples planned deficits in its target bundles to the
mass that will supply them.  One planned count unit subtracts mass $g$
from each target bundle and adds mass $g$ at the supplier if it lies in $U$.
A bundle supplier incurs no additional opening in this accounting.
Thus a star with $d$ targets changes the opening mass by $(1-d)g$ for
an outside supplier and by $-dg$ for a bundle supplier.

Make a list $\mathcal E$ of stars; a star $e$ has weight $w_e$, a flag
$a_e\in\{0,1\}$ indicating an outside supplier, a target set
$D_e\subseteq\mathcal S$, and a specified supplier:
\begin{enumerate}
 \item For each $i\in U$, take $w_e=y_i$, $a_e=1$,
 $D_e=\{s:i\in A_s\}$, and supplier $i$.
 \item For each pair $s,t$ with $y(A_s\cap P_t)>0$, take
 $w_e=y(A_s\cap P_t)$, $a_e=0$, $D_e=\{s\}$, and supplier $P_t$.
 Necessarily $t\ne s$.
 \item For each $s$ with $c_s>0$, take $w_e=\gamma c_s$, $a_e=0$,
 $D_e=\{s\}$, and use as supplier the bundle of a nearest other
 representative.  Such a representative exists because $c_s>0$.
\end{enumerate}
The last stars are called extra stars.  The planned count of star $e$
is $q_e=w_e/g$.  With $u_e=a_e-|D_e|$, the planned raw opening mass is
\begin{equation}\label{eq:prep-raw-identity}
 |\mathcal S|+\sum_e u_e w_e
 =\sum_i y_i-\gamma\sum_s c_s.
\end{equation}
To see this, the flag-one weights sum to $y(U)$, while the total weight
targeting $s$ is $(1+\gamma)c_s$.

Apply Lemma~\ref{lem:prep-weighted-rounding} independently to the
positive and negative signs of $u_e$, using weights $|u_e|$.  Complete
the rounding of the positive class by rounding its possible last
fraction with the correct mean.  Round the zero-coefficient counts
independently between their floor and ceiling.  Leave the last
fractional count of the negative class, if any, temporarily unrounded.
All counts are then integral except possibly this one.

If there is no remaining fraction, the resulting counts are both the
nominal counts and the actual activation counts.  Otherwise, write the
remaining count as $m+\xi$, where $0<\xi<1$, and let its star have
$d=|D_e|$ targets.  Its nominal count is $m+\operatorname{Bernoulli}(\xi)$.
If its flag is one, open $(m+1)g$ at its supplier, irrespective of the
nominal coin.  Activate $m$ units at each target, and activate one
additional unit at every target in a random subset of size
\[
 H\in\{\lfloor d\xi\rfloor,\lceil d\xi\rceil\},
 \qquad \E H=d\xi,
\]
whose membership marginal is $\xi$ at every target.  Such a subset is
obtained by choosing the indicated size with the correct mean and
then taking a uniformly shifted interval of that size in a fixed
cyclic ordering of the targets.  The nominal coin can be independent
of this choice.  All other stars use their integral counts for both
openings and activations.

This modification bounds the raw opening error independently of the
number of targets.  For the exceptional star its actual signed
contribution, in count units, minus its planned contribution is
\[
 a_e(m+1)-dm-H-(a_e-d)(m+\xi)
 =a_e(1-\xi)+d\xi-H<2.
\]
A positive coefficient $u_e$ necessarily equals one, so the last
rounding in the positive class adds less than one further count unit.
The zero class has no signed contribution.  If $\widetilde y_i$ denotes
the actual opening on $U$, and $Z_s$ denotes the total activated deficit
mass at $s$, we therefore have on every outcome
\begin{equation}\label{eq:prep-raw-bound}
 \sum_{i\in U}\widetilde y_i+\sum_s(1-Z_s)
 \le\sum_i y_i-\gamma\sum_s c_s+3g.
\end{equation}
Moreover, each actual outside opening dominates both its nominal
opening and every activation demand placed on it by a single target.

\paragraph{Deficit moments and clipping.}
The signed count already controls raw opening mass. The remaining
obstacle is that an activated deficit might leave too little mass in its
own bundle. We cap that deficit and pay for the added bundle mass with
the reserve in Equation~\ref{eq:prep-raw-bound}.
For every $s$, let $A_{es}$ be the integer activation count from a star
targeting $s$.  Then
\begin{equation}\label{eq:prep-deficit-moments}
 \E Z_s=(1+\gamma)c_s,
 \qquad \Var(Z_s)\le g(1+\gamma)c_s.
\end{equation}
Here it matters that the last negative star may depend on the preceding
random choices.  Condition on the complete count vector immediately
before its exceptional rounding.  At most one coordinate is
fractional, so for distinct $e,f$ targeting $s$, at least one of
$A_{es},A_{fs}$ is deterministic under this conditioning.  Each has
conditional mean equal to its current count.  Their conditional product
expectation is therefore the product of those two current counts.
The distinct-pair product bound in
Lemma~\ref{lem:prep-weighted-rounding} consequently still holds for
these activations.  Their individual means are $q_e$, and each is
supported on $\{\lfloor q_e\rfloor,\lceil q_e\rceil\}$.  Thus
\[
 \Var(Z_s)
 \le g^2\sum_{e:s\in D_e}\Var(A_{es})
 \le g^2\sum_{e:s\in D_e}q_e
 =g(1+\gamma)c_s,
\]
as asserted.  The nominal counts also retain the same first and
distinct-pair moment bounds.

Clip the deficit at one half:
\[
 Z_s^*=\min(Z_s,\tfrac12),\qquad p'_s=1-Z_s^*.
\]
Discard any excess activation units.  This can be done in whole units
because $\Delta$ is even.  By Equation~\ref{eq:prep-parameter-basics},
$\mu_s:=\E Z_s\le1/4$.  For any random variable $X$ with mean $\mu<b$,
\[
 (X-b)_+\le\frac{(X-\mu)^2}{b-\mu}.
\]
Indeed the assertion is immediate when $X\le b$, and otherwise follows
by expanding $(X-\mu)^2=(X-b+b-\mu)^2$.
Equation~\ref{eq:prep-deficit-moments} therefore gives
\[
 \E(Z_s-\tfrac12)_+\le4g(1+\gamma)c_s.
\]
Let $\mathcal G$ be the event that
\[
 \sum_s(Z_s-Z_s^*)\le\gamma\sum_s c_s.
\]
If some $c_s$ is positive, Markov's inequality gives
\begin{equation}\label{eq:prep-budget-probability}
 \Prob(\mathcal G)\ge1-\rho,
 \qquad \rho:=\frac{4g(1+\gamma)}\gamma.
\end{equation}
If all $c_s$ vanish, all activations and discarded amounts are zero,
and $\mathcal G$ holds surely.  Equations~\ref{eq:prep-raw-bound} and
\ref{eq:prep-budget-probability} show that a vector with bundle sums
$p'_s$ has total opening mass at most $\sum_i y_i+3g$ on $\mathcal G$.

\paragraph{Rounding within bundles and supplying one unit.}
Condition on the complete star-rounding outcome, including the retained
activations.  In each bundle independently, rescale piece $i\in P_s$ to
mass
\[
 \overline y_i=y_i\frac{p'_s}{p_s}.
\]
Since $y_i\le p_s$ and $1/2\le p'_s\le1$, these masses belong to $[0,1]$.
Their sum $p'_s$ is a grid multiple.  Apply the unweighted version of
Lemma~\ref{lem:prep-weighted-rounding} to counts $\overline y_i/g$.
The preserved sum is integral in count units, so there is no final
fraction.  This produces grid masses $\widetilde y_i$ with conditional
means $\overline y_i$ and nonpositive conditional pair covariances.
All masses remain at most one, because one is a grid point.  Likewise
the outside opening at a piece of original mass $y_i$ is at most
$g\lceil y_i/g\rceil\le1$, even under the exceptional correction.

Every representative $s$ can now receive one unit: take its entire
bundle mass $p'_s$ and supply its retained deficit $Z_s^*$ using the
retained activations.  Each outside supplier has enough mass for its
activation demand at $s$.  Any bundle supplier is different from $P_s$.
The total demand of $s$ on any one such bundle is at most its entire
retained deficit, hence at most $1/2$, while every supplying bundle has
mass at least $1/2$.  Allocate this demand greedily over its pieces.
The bundles and $U$ are disjoint, and different clients have separate
assignment constraints.  This proves feasibility for each client
individually, even before conditioning on $\mathcal G$.  It also shows
that the total opening mass is at least one.

\paragraph{Cost comparison at a distant representative scale.}
Fix a client $j$ and its representative $s$, and put $D=d(j,s)$.
Let $R_s=\max_{i\in F_s}d(s,i)$, and define the radius
\[
 K_s=\begin{cases}L_s,&c_s>0,\\R_s,&c_s=0.\end{cases}
\]
If $V_j=0$, Equation~\ref{eq:prep-representative-properties} gives
$D=V_s=0$.  All pieces of $F_s$ then have distance zero from $j$,
$c_s=0$, and $Z_s=0$ surely.  The rounded bundle $P_s=F_s$ still has
one unit at distance zero.  Hence the new cost is identically zero.
We may therefore assume $V_j>0$ in the remaining cost analysis.

Suppose first that $K_s>T_1V_j$.  We have $R_s\ge\delta K_s$: this is
immediate if $c_s=0$, and otherwise some assigned mass lies beyond
$\delta L_s$.  Consequently $D/R_s\le\alpha<1$.  This controls the
cost of exchanging the portions of $F_j$ and $F_s$ that are not shared.
Leave the shared pieces matched identically.  The masses of
$F_s\setminus F_j$ and $F_j\setminus F_s$ are equal.  Each piece in the
former set is at distance at most $R_s+D$ from $j$.  Each piece in the
latter set is at distance at least $R_s-D$ from $j$, because the
nearest-unit assignment at $s$ uses all closer mass before stopping at
radius $R_s$.  This argument remains valid at ties.  Comparing the two
residual costs shows that using the old assignment $F_s$ for client $j$
has cost at most
\begin{equation}\label{eq:prep-far-exchange}
 C_j(F_s):=\sum_{i\in F_s}y_i d(j,i)
 \le\frac{1+\alpha}{1-\alpha}V_j.
\end{equation}

Serve $j$ with the feasible assignment just constructed for $s$.
For each $i\in P_s$,
\[
 \E\widetilde y_i
 =y_i\E[p'_s]/p_s\le y_i/p_s.
\]
For an ordinary star, the expected retained activation mass is at most
its original weight.  Thus outside suppliers do not increase their
expected cost charge.  It remains to compare arbitrary new supply
inside a bundle $P_t$ with the old pieces in $A_s\cap P_t$.
This case only arises when $c_s>0$.  Then $D\le\beta L_s$, and for any
old supplier piece $i\in P_t$,
\[
 d(j,i)\ge d(s,t)-D-\delta L_t
 \ge(1-\beta-\delta)d(s,t).
\]
Any two pieces of $P_t$ are at distance at most $2\delta L_t$.
Because $L_t\le d(s,t)$ and $1-\beta-\delta>0$, every possible new
supplier piece has distance to $j$ at most
\[
 \kappa\,d(j,i),\qquad
 \kappa:=1+\frac{2\delta}{1-\beta-\delta}.
\]
The comparison holds for every old piece, so it also holds relative to
the weighted mean distance of the old pieces aggregated by that star.
Multiplying by its expected retained activation mass bounds its
expected cost by $\kappa$ times its old cost contribution.

Finally, for the extra star at $s$, its supplier $P_t$ satisfies
$d(s,t)=L_s$ and $L_t\le L_s$.  Every supply position is therefore at
distance at most $(\beta+1+\delta)L_s$ from $j$.  By
Equation~\ref{eq:prep-deficit-small}, its expected cost is at most
\[
 (\beta+1+\delta)\gamma c_s L_s
 \le\frac{(\beta+1+\delta)\gamma}{\delta}V_j.
\]
There is no extra star when $c_s=0$.
Combining these bounds with Equation~\ref{eq:prep-far-exchange}, the
unconditional expected nearest-unit cost in this case is at most
$B_{\mathrm{far}}V_j$, where
\begin{equation}\label{eq:prep-far-factor}
 B_{\mathrm{far}}
 :=\max\{p_0^{-1},\kappa\}\frac{1+\alpha}{1-\alpha}
   +\frac{(\beta+1+\delta)\gamma}{\delta}.
\end{equation}

\paragraph{Cost comparison at a nearby representative scale.}
Now suppose $K_s\le T_1V_j$.  Consider the mass on the old set $F_j$,
using nominal rounded masses on $U$ and actual rounded masses in the
bundles.  This mass is available in the actual opening vector.  Its
unconditional expected cost is at most $V_j/p_0$: nominal outside
openings have mean $y_i$, and a bundle piece has expected mass at most
$y_i/p_s\le y_i/p_0$.

We bound the expected shortfall of this available mass from one.
Write $M_U$ for its nominal mass on $U\cap F_j$.  The first and pair
moments of the nominal counts imply
\begin{equation}\label{eq:prep-outside-noise}
 \E M_U=y(U\cap F_j),\qquad
 \Var(M_U)\le g\,y(U\cap F_j)\le g.
\end{equation}
For a fixed complete star outcome, put
\[
 a_t=y(F_j\cap P_t),\qquad
 \overline M_P=\sum_t a_t p'_t/p_t,
\]
and write $M_P$ for the actual rounded mass on these bundle pieces.
Conditional on the star outcome, bundle rounding has nonpositive pair
covariances within each bundle and is independent between bundles.
The variance of a single grid-rounded mass of conditional mean
$\overline y_i$ is at most $g\overline y_i$.  Therefore
\begin{equation}\label{eq:prep-conditional-noise}
 \E[M_P\mid\text{star outcome}]=\overline M_P,\qquad
 \Var(M_P\mid\text{star outcome})
 \le g\overline M_P\le g/p_0.
\end{equation}

Clipping at $1/2$ cannot increase distance from $c_t\le1/2$, so
Equation~\ref{eq:prep-deficit-moments} gives
\begin{align}
 \E|p'_t-p_t|
 &=\E|c_t-Z_t^*|\notag\\
 &\le\E|c_t-Z_t|
 \le\gamma c_t+\sqrt{g(1+\gamma)c_t}.
 \label{eq:prep-rescaling-noise}
\end{align}
Since $\sum_t a_t\le1$, it follows that
\[
 \E\left|\overline M_P-\sum_t a_t\right|
 \le\sum_t\frac{a_t}{p_t}\E|p'_t-p_t|
 \le\frac{\gamma+\sqrt{g(1+\gamma)}}{p_0}.
\]
Use the identity $y(U\cap F_j)+\sum_t a_t=1$, the triangle inequality,
and the second-moment bounds in
Equations~\ref{eq:prep-outside-noise}--\ref{eq:prep-conditional-noise}.
The expected shortfall is at most
\begin{equation}\label{eq:prep-shortfall}
 \E(1-M_U-M_P)_+\le H,
 \qquad
 H:=\sqrt g+\sqrt{g/p_0}
       +\frac{\gamma+\sqrt{g(1+\gamma)}}{p_0}.
\end{equation}
In particular no independence between outside counts and rescaled
bundle masses is being assumed.

Trim any surplus mass on $F_j$, which cannot increase its cost.  To
fill a remaining shortfall, there is always one unit of available
opening within distance
\[
 QV_j,\qquad Q:=M+(1+\delta)T_1,
\]
from $j$.  Indeed, if $c_s=0$, there are no activations at $s$ and $P_s$
has exactly one unit, all at distance at most
$D+R_s\le(M+T_1)V_j$.  If $c_s>0$, use $P_s$ and the bundle $P_t$ of a
nearest other representative.  They are disjoint and each has at
least half a unit.  Their supply lies within distance
$D+(1+\delta)L_s\le QV_j$.
This supply can fill the shortfall even if some of it is already used:
if the current assignment has mass $1-q$, a ball containing at least
one opening unit contains at least $q$ unused units.  Allocate those
units to finish the assignment.  By Equation~\ref{eq:prep-shortfall},
the unconditional expected nearest-unit cost is at most
$B_{\mathrm{near}}V_j$, where
\begin{equation}\label{eq:prep-near-factor}
 B_{\mathrm{near}}:=p_0^{-1}+QH.
\end{equation}

\paragraph{Choosing parameters and enforcing the budget.}
Equations~\ref{eq:prep-far-factor} and~\ref{eq:prep-near-factor} give
explicit bounds that approach one in the stated order of parameter
choices.  More precisely, first choose $\delta$ small to make the
limiting value of $\kappa$ close to one.  Next choose $M$ large so that
$p_0^{-1}$ is close to one, and then choose $T_1$ large so that $\alpha$
and $\beta$ are small.  With these constants fixed, choose $\gamma$
small enough to control both the last term of $B_{\mathrm{far}}$ and
the term $Q\gamma/p_0$ in $B_{\mathrm{near}}$.  Finally choose $g$ small
enough to control all square-root terms and $\rho$ in
Equation~\ref{eq:prep-budget-probability}.  All constants may be
rational, and all sufficiently small $g=1/\Delta$ with even $\Delta$
satisfy the resulting strict bounds.  In particular, arrange
\[
 \rho<\tfrac12,\qquad
 \frac{\max\{B_{\mathrm{far}},B_{\mathrm{near}}\}}{1-\rho}
 \le1+\eta.
\]

Reject a star-rounding outcome outside $\mathcal G$ and repeat.
For the purpose of the preceding unconditional analysis, the subsequent
bundle rounding can be imagined on every outcome.  Let
$\widetilde V_j$ denote the resulting nearest-unit cost.  Its
nonnegativity, rather than any independence from $\mathcal G$, gives
\begin{equation}\label{eq:prep-cost-final}
 \E[\widetilde V_j\mid\mathcal G]
 \le\frac{\E\widetilde V_j}{\Prob(\mathcal G)}
 \le(1+\eta)V_j.
\end{equation}
The already established supply and opening bounds hold on every
accepted outcome.

The split representation, star list, and number of pair updates are
polynomial in the input size.  Pair updates and subset choices use
rational probabilities and rational linear operations.  Their
encodings have polynomial bit length: each weighted update introduces
only ratios of integer weights bounded by $|\mathcal S|$, and there
are polynomially many updates; the subsequent rescalings divide by
the rational masses $p_s$.  Exact rational sampling has expected
polynomial running time.  The acceptance probability is at least
one half, so rejection multiplies the expected running time by at
most two.  This proves the theorem.
\end{proof}
